\section*{Bab \ref{ch:b2:plant-plasticity} --- Adaptasi Tumbuhan dan Keplastisan Fenotipe}

\begin{solution}{exo:b2:plant-plasticity:1}
Keplastisan: satu \omterm{def:b2:meiosis-heredity:vocabulary}{genotipe}, beberapa \omterm{def:b2:meiosis-heredity:vocabulary}{fenotipe} menurut lingkungannya
(\omterm{prop:b2:plant-plasticity:sunshade}{daun matahari} dan \omterm{prop:b2:plant-plasticity:sunshade}{daun naungan} pada satu pohon ek). \omterm{def:b2:plant-plasticity:plasticity}{Norma reaksi}:
kurva \omterm{def:b2:meiosis-heredity:vocabulary}{fenotipe} sebuah \omterm{def:b2:meiosis-heredity:vocabulary}{genotipe} terhadap sebuah peubah lingkungan
(ketebalan daun terhadap cahaya). \omterm{def:b2:plant-plasticity:plasticity}{Aklimasi}: penyesuaian faali yang
dapat balik (pengerasan gandum hitam terhadap embun beku pada musim
gugur). \omterm{def:b2:plant-plasticity:plasticity}{Adaptasi}: kecocokan terwariskan yang terseleksi selama banyak
generasi (metabolisme CAM sebatang kaktus).
\end{solution}

\begin{solution}{exo:b2:plant-plasticity:2}
Koleoptil utuh yang disinari dari satu sisi: membengkok menuju
cahayanya. Ujung yang ditudungi tudung tak tembus cahaya: tanpa
pembengkokan --- ujungnya mengindra cahayanya. Ujung yang ditudungi
tudung bening: membengkok --- tudungnya sendiri tak berbahaya. Pangkal
yang dilindungi tabung tak tembus cahaya dengan ujung terbuka:
membengkok --- pengindraannya di ujungnya dan tanggapannya di bawahnya.
\end{solution}

\begin{solution}{exo:b2:plant-plasticity:3}
\omterm{prop:b2:plant-plasticity:sunshade}{Daun matahari}: kecil, tebal, dua atau tiga lapis sel tiang, kutikula
tebal, banyak stoma, fotosintesis maksimal yang tinggi, dan titik
kompensasi yang tinggi. \omterm{prop:b2:plant-plasticity:sunshade}{Daun naungan}: lebar, tipis, satu lapis sel
tiang, lebih banyak klorofil tiap gram, respirasi yang rendah, titik
kompensasi yang rendah. Pilihannya diambil ketika daunnya masih berupa
primordium di dalam tunasnya, dari cahaya yang ia terima; sehelai daun
yang masak tidak dapat berubah.
\end{solution}

\begin{solution}{exo:b2:plant-plasticity:4}
\omterm{prop:b2:plant-plasticity:stomata}{Sel penjaganya} menyerap kalium dan air, menggembung, dan, karena
dinding dalamnya lebih tebal, melengkung saling menjauh untuk membuka
porinya. Cahaya biru dan $\mathrm{CO_2}$ dalaman yang rendah membukanya
pada pagi hari; \omterm{def:b2:plant-flowering:dormancy}{asam absisat}, dari akar di dalam tanah yang mengering
atau dari daunnya, menutupnya pada sebuah kekeringan, begitu pula udara
yang sangat kering.
\end{solution}

\begin{solution}{exo:b2:plant-plasticity:5}
$\varphi = R/(R + 0.77\,FR)$ dengan $FR = 1$: $1.15$: $0.60$; $0.7$:
$0.48$; $0.2$: $0.21$; $0.05$: $0.06$. Ambangnya $0.4$ jatuh di antara
rumpangnya ($0.48$) dan kanopinya ($0.21$): \omterm{thm:b2:plant-plasticity:phytochrome}{penghindaran naungannya}
menyala di bawah kanopinya.
\end{solution}

\begin{solution}{exo:b2:plant-plasticity:6}
$E = 0.25\times 0.025 = \qty{6.25}{mmol.m^{-2}.s^{-1}}$; $A =
0.25\times 100/1.6 = \qty{15.6}{\micro mol.m^{-2}.s^{-1}}$; $A/E =
2.5\times 10^{-3}$ mol karbon tiap mol air, yakni 400 molekul air tiap
karbon, $400\times 18/12 = \qty{600}{g}$ air tiap gram karbon. Dalam 12
jam: $6.25\times 10^{-3}\times 43200 =
\qty{270}{mol/m^2}$,
\qty{4.9}{L} tiap meter persegi.
\end{solution}

\begin{solution}{exo:b2:plant-plasticity:7}
$E = 0.25\times 0.005 = \qty{1.25}{mmol.m^{-2}.s^{-1}}$, $A$ tetap:
$A/E = 0.0125$, 80 air tiap karbon, \qty{120}{g/g} --- seperlima ongkos
siang harinya.
\end{solution}

\begin{solution}{exo:b2:plant-plasticity:8}
$v = 2r^{2}\Delta\rho g/9\eta = 2\times 4\times 10^{-12}\times
400\times 9.81/0.09 = \qty{3.5e-7}{m/s}$: \qty{0.35}{\micro m/s},
kira-kira \qty{30}{s} untuk melintasi selnya. Pada sebuah klinostat
arah gravitasi terhadap selnya berubah tiap beberapa detik pengendapan,
sehingga butirnya tidak pernah menumpuk pada satu sisi dan rangsangan
reratanya selama waktu pemaduan tumbuhannya (menit) adalah nol.
\end{solution}

\begin{solution}{exo:b2:plant-plasticity:9}
Pemanjangan rerata dalam dua jam $0.04$; sisinya $1.2\times 0.04 = 0.048$ dan $0.8\times 0.04 = 0.032$; $\theta = 20\times 0.016/2 = 0.16$
rad, kira-kira \qty{9}{\degree}. Pembengkokannya menumpuk sebesar
$0.08$ rad tiap jam; $\pi/2$ memakan kira-kira \qty{20}{h}.
\end{solution}

\begin{solution}{exo:b2:plant-plasticity:10}
Cadangannya memungkinkan batang \qty{50}{cm}. Rumpun jelatangnya:
\qty{30}{cm} dalam 7.5 hari dengan \qty{30}{mg} --- semainya mencapai
cahaya dengan cadangan yang tersisa. Hutannya: \qty{10}{m} tak
terjangkau; semainya membelanjakan \qty{50}{mg}-nya pada setengah meter
batang yang pucat dalam dua belas hari lalu mati teretiolasi, sedangkan
siasat yang tahan naungan --- beberapa daun tipis, pertumbuhan yang
lambat --- mungkin dapat menjaganya tetap hidup selama bertahun-tahun.
\end{solution}

\begin{solution}{exo:b2:plant-plasticity:11}
Es terbentuk lebih dahulu di ruang antarselnya, tempat larutannya lebih
encer; tekanan uap di atas es lebih rendah daripada di atas air selnya,
sehingga airnya meninggalkan selnya untuk bergabung dengan esnya dan
selnya mengalami dehidrasi lalu membrannya runtuh. Sepekan dingin
memuati selnya dengan gula dan protein pelindung yang menurunkan titik
bekunya dan memantapkan membrannya, serta membuat protein antibeku yang
membatasi pertumbuhan hablurnya. Menyalakan programnya sepanjang tahun
mengalihkan karbon dari pertumbuhan ke perlindungan dan memperlambat
pembelahan: tumbuhan yang mengeras secara tetap adalah tumbuhan yang
kecil.
\end{solution}

\begin{solution}{exo:b2:plant-plasticity:12}
Isyarat merah jauhnya jujur ketika tetangganya dapat dilampaui dan
dusta ketika tetangganya sebatang pohon; sepekan yang hangat jujur pada
bulan April dan dusta pada bulan Februari; di dalam tanaman yang rapat
tetangga tiap tumbuhan adalah sepadannya, sehingga isyaratnya jujur
dalam bentuk tetapi tak berguna, dan pemulia menyeleksi tumbuhan yang
mengabaikannya. Sebuah \omterm{def:b2:plant-plasticity:plasticity}{norma reaksi} adalah catatan seberapa sering tiap
isyarat telah mengatakan yang sebenarnya pada masa lalu garis
keturunannya, dan ia gagal setiap kali masa kini berbeda dari masa lalu
itu.
\end{solution}

\begin{solution}{pb:b2:plant-plasticity:1}
\textbf{1.} Terbuka: $1.15/(1.15 + 0.77) = 0.60$; kanopi: $0.2/(0.2 +
0.77) = 0.21$.
\textbf{2.} Nilai $\varphi$ yang rendah: ada tumbuhan di atasnya ---
memanjang. Di bawah sebuah kain naung $\varphi$ tetap $0.60$: sebuah
hari yang gelap, bukan seorang tetangga; tanpa \omterm{thm:b2:plant-plasticity:phytochrome}{penghindaran naungan},
hanya pertumbuhan yang lebih lambat karena kurang cahaya.
\textbf{3.} $1 + 2(0.60 - 0.21) = 1.78$: \qty{2.7}{cm} sehari.
\textbf{4.} \qty{37}{cm} berbanding \qty{21}{cm}.
\textbf{5.} Pembengkokan menuju rumpangnya (fototropisme, lewat
reseptor cahaya biru fototropin) serta pemanjangan yang lebih cepat dan
penataan ulang daun pada sisi yang $R{:}FR$-nya lebih tinggi (lewat
fitokrom).
\textbf{6.} \omterm{prop:b2:plant-plasticity:sunshade}{Daun naungan} dibangun lebar dan tipis untuk menangkap
cahaya yang langka dengan murah; di bawah matahari penuh daun itu
kepanasan, terhambat cahaya, lalu hangus, dan tumbuhannya harus
menggantinya dengan \omterm{prop:b2:plant-plasticity:sunshade}{daun matahari}.
\textbf{7.} \qty{2.7}{cm} sehari sepanjang zona \qty{15}{mm} adalah
\qty{0.11}{cm} tiap jam, yaitu pemanjangan nisbi $0.074$ tiap jam;
sisi yang ternaung $1.3\times 0.074 = 0.096$, sisi yang tersinari $0.7\times 0.074 =
0.052$.
\textbf{8.} $\theta = 15\times(0.096 - 0.052)/2 = 0.33$ rad, kira-kira
\qty{19}{\degree} dalam sejam.
\textbf{9.} \qty{60}{\degree} adalah $1.05$ rad: kira-kira tiga jam.
\textbf{10.} $v = 2\times 6.25\times 10^{-12}\times 400\times
9.81/0.09 = \qty{5.5e-7}{m/s}$, \qty{0.55}{\micro m/s}: \qty{22}{s}
untuk melintasi selnya.
\textbf{11.} Fototropisme, dan \omterm{thm:b2:plant-plasticity:phytochrome}{penghindaran naungan} itu sendiri
menaikkan sudut yang membuat pucuknya puas untuk tumbuh: cahayalah
sumber daya yang diperebutkan, dan batang yang meluruskan diri melawan
gravitasi akan berbalik masuk ke dalam naungan.
\textbf{12.} Akarnya membengkok ke bawah sampai ia tegak: statolitnya
mengendap ke sisi bawahnya, auksin menumpuk di sana, dan di dalam
sebuah akar kelebihan auksin itu menghambat pertumbuhan, sehingga sisi
atasnya tumbuh lebih cepat. Pengalihan yang sama, kepekaan sel yang
berlawanan.
\textbf{13.} $E = 0.2\times 0.015 = \qty{3}{mmol.m^{-2}.s^{-1}}$;
sepanjang $2\times 10^{-3}$ \unit{m^2} dan \qty{50400}{s}:
\qty{0.30}{mol}, \qty{5.4}{g} air sehari.
\textbf{14.} $A = 0.2\times 80/1.6 = \qty{10}{\micro mol.m^{-2}.s^{-1}}$;
tiap hari $0.50$ mol karbon tiap meter persegi; $A/E = 3.3\times
10^{-3}$: 300 air tiap karbon, \qty{450}{g/g}.
\textbf{15.} $0.50\times 2\times 10^{-3} = \qty{1}{mmol}$: \qty{12}{mg}
karbon sehari.
\textbf{16.} Rumpang: $E = \qty{6}{mmol.m^{-2}.s^{-1}}$, \qty{10.9}{g}
sehari; daya gunanya \qty{900}{g/g}.
\textbf{17.} $E = 0.02\times 0.03 = \qty{0.6}{mmol.m^{-2}.s^{-1}}$
(turun sepuluh kali lipat); $A = 0.02\times 200/1.6 = \qty{2.5}{\micro
mol.m^{-2}.s^{-1}}$ (turun empat kali lipat). Airnya turun lebih
banyak: ketika daunnya menurunkan $\mathrm{CO_2}$ dalamannya landaian
bagi karbonnya melebar sedangkan landaian bagi airnya tidak dapat.
\textbf{18.} \qty{4.25}{g} air, \qty{0.85}{g} yang boleh hilang; pada
\qty{10.9}{g} sehari ($0.78$ g tiap jam) ia layu dalam kira-kira sejam.
\textbf{19.} CAM lambat dan memerlukan jaringan penyimpan yang sukulen;
sebuah semai yang berpacu mengejar cahaya di bawah sebuah kanopi,
tempat kehilangan airnya kecil, memerlukan kecepatan, bukan penghematan.
\textbf{20.} $100/2 = \qty{50}{cm}$.
\textbf{21.} Rumpun jelatang: \qty{40}{cm} dalam 15 hari dengan
\qty{80}{mg} --- ia mencapai cahaya. Hutan: \qty{8}{m} tak terjangkau;
ia membelanjakan cadangannya pada setengah meter lalu mati. Yang lebih
baik di dalam hutan: tetap pendek, membuat \omterm{prop:b2:plant-plasticity:sunshade}{daun naungan} yang tipis, dan
menunggu sebuah rumpang --- caranya beech.
\textbf{22.} Birch: norma yang curam, memanjang dengan kuat pada
$\varphi$ yang rendah, karena tetangga seorang perintis sebesar dirinya
sendiri. Beech: norma yang datar, sedikit pemanjangan dan banyak
ketahanan, karena tetangganya adalah pepohonan.
\textbf{23.} Lumpuhkan tanggapan \omterm{thm:b2:plant-plasticity:phytochrome}{penghindaran naungannya} --- jaga agar
fitokrom B tetap aktif atau singkirkan \omterm{prop:b2:cell-differentiation:transcriptional}{faktor transkripsi} yang
ditahannya --- lalu periksa waktu pembungaan, perkecambahan, dan
kekuatan batangnya, yang dikendalikan lintasan yang sama.
\textbf{24.} Tumbuhan yang tidak dapat mengindra tetangganya terlampaui
setiap kali ia sebenarnya dapat menang; tumbuhan yang selalu memercayai
isyaratnya menyia-nyiakan dirinya di bawah tiap pohon. Norma yang
menang bertaruh pada isyaratnya sebanding dengan seberapa sering, pada
masa lalu garis keturunannya, isyaratnya benar.
\textbf{25.} $\varphi = 0.60$ di tempat terbuka dan $0.21$ di bawah
kanopinya; \qty{37}{cm} setelah 14 hari; \qty{5.4}{g} air sehari di
dalam naungannya, \qty{10.9}{g} di dalam rumpangnya.
\end{solution}
